\chapter{绪论}

\section{研究背景与意义}

时空序列预测是指利用历史时刻的观测值及其所附带的时间戳、空间位置或拓扑关系等信息，预测未来一段时间内多个变量或空间节点取值的数据挖掘任务。它既是时间序列预测与图数据学习的交叉问题，也是智慧城市、智能交通、能源管理、气象预报、工业监测和金融风控等实际应用中的基础性任务\cite{traffic-stf-survey,gnn4ts,survey01}。随着物联网、移动传感和城市数字化程度的不断提高，各类业务系统持续产生海量、高频、高维且高度耦合的时空数据。对这些数据进行准确预测，不仅有助于提前识别交通拥堵、设备异常和供需波动等风险，也能为资源调度、运营决策和风险控制提供重要依据，因而具有突出的理论价值与现实意义。

按照数据中是否显式给定空间图结构，时空序列预测可以大致划分为两类典型问题。一类是面向一般多元时间序列的预测任务，其输入通常为多个通道同步采集的时间序列，如电力负荷、汇率、气象要素等，变量之间可能存在隐含的相关关系；另一类是以交通流预测为代表的空间图时间序列预测任务，其输入通常对应交通路网中的若干传感器或路段节点，节点之间通过预先给定或可学习的邻接矩阵表达空间连接关系。两类任务在输入形式上有所区别，但都面临着若干共同挑战：观测数据往往呈现明显的非平稳性，统计特性随时间或空间尺度发生变化；序列中同时存在趋势、周期、局部波动等多尺度模式；不同通道或不同区域的数据分布存在显著差异；时空依赖既包含由物理连接决定的静态结构，也包含随时间演化的动态关系。

传统时间序列预测方法，如自回归移动平均模型、指数平滑和季节趋势分解等，基于较为明确的统计假设对序列进行建模，在平稳、低维和线性场景中具有良好的可解释性\cite{brockwell1991time,1990STL}。然而，现实时空序列通常具有非线性、长程依赖、高维耦合和非平稳等特点，使得传统方法的表达能力受到明显限制。近年来，以循环神经网络、卷积神经网络、图神经网络和 Transformer 为代表的深度学习方法被广泛应用于时空序列预测，并通过注意力机制、频域分析、实例归一化和多尺度分解等手段显著提升了预测精度\cite{wen2023transformers,gnn4ts,kim2024comprehensive}。尽管如此，现有方法在刻画时空序列内在异质性与尺度变化方面仍存在不足：多数模型采用单一尺度的特征提取方式，或对所有变量、所有节点共享同一组参数，难以自适应地处理不同尺度和不同分布模式；部分方法虽然引入多尺度结构，但尺度划分通常依赖经验或先验知识，无法根据当前输入动态调整；同时，对时间戳信息、静态拓扑与动态时空依赖的联合利用仍不充分。

混合专家（Mixture-of-Experts，MoE）机制为上述问题提供了一种自然的建模思路。其核心思想是针对输入样本或特征的分布差异，将其路由到若干具有专门化能力的专家网络中进行处理，再对专家输出进行加权融合，从而在不显著增加整体计算量的前提下提升模型对异构模式的适应能力\cite{moe,moe-survey}。时空序列中的非平稳性和异质性恰好与 MoE 的“分而治之”思想相契合：不同时间尺度、不同变量通道以及不同空间区域可以对应不同的专家，从而避免共享参数模型在分布差异较大时的表达能力退化。然而，如何针对时空序列设计合适的路由信号、专家结构、多尺度信息融合方式，以及如何同时建模静态与动态依赖，仍然是尚未充分解决的问题。

基于上述分析，本文围绕“多尺度建模”和“混合专家系统”两条主线展开研究，分别面向一般时间序列预测和时空交通预测任务提出两类模型框架。前者聚焦时间维度的双向时间戳依赖、多尺度时序模式和通道间分布差异，构建多尺度时序混合专家与双向时间戳依赖建模框架；后者聚焦空间维度的多尺度图结构以及静态图结构与动态时空依赖共存问题，构建多尺度空间自适应建模与时空静动态混合专家框架。两项工作从时间域与空间域两个侧面共同验证了“多尺度建模 + 混合专家”在时空序列预测中的有效性，为构建更具适应性和鲁棒性的时空预测模型提供了新的技术路线。

\section{国内外研究现状}

\subsection{基于深度学习的多元时间序列预测方法}

时间序列预测旨在根据历史观测对未来若干时刻的取值进行推断，是时序数据挖掘中的基础任务\cite{kim2024comprehensive,wen2023transformers}。早期深度学习方法主要沿三类网络结构展开。以循环神经网络为基础的方法利用序列递归建模捕捉时间依赖，例如 DeepAR 将自回归思想与循环网络相结合，用于短期序列预测\cite{deepar}。以卷积神经网络为基础的方法通过时间卷积和局部感受野提取时序特征，具有训练高效、易于并行化的优点\cite{micn,cnn_related1}。以图神经网络为基础的方法则把变量映射为图节点，通过学习变量之间的隐式关系辅助预测，适用于变量之间存在显著相关性的多元序列\cite{cdotsgnn,mtgnn,gnnsuv}。

Transformer 的兴起显著推动了长程时序建模的发展。Informer 通过稀疏注意力降低长序列建模的计算复杂度\cite{informer}；Autoformer 引入序列分解与自相关机制，以显式刻画趋势和周期成分\cite{autoformer}；PatchTST 将时间序列切分为若干时间块并建模块间依赖\cite{patchtst}；Crossformer 同时关注时间维度和变量维度的跨维依赖\cite{crossformer}；iTransformer 则把变量作为令牌进行处理，使注意力机制更适合建模变量间关系\cite{itransformer}。与此同时，DLinear 等线性模型表明，在部分时序预测基准上，简单线性结构同样具有较强的竞争力，促使研究者重新审视模型归纳偏置与任务特性之间的匹配\cite{dlinear}。

针对时间序列中的多尺度模式，Pyraformer 采用金字塔注意力在多个时间尺度上建模长程依赖\cite{pyraformer}；SCINet 通过二分下采样树进行层次化分析\cite{scinet}；TimeMixer 在多个尺度上分别分解并融合趋势项和季节项\cite{timemixer}；Pathformer 将输入自适应划分为不同长度的分块，以提取尺度特异性特征\cite{pathformer}；PDF 则从周期解耦的角度对长时预测进行建模\cite{pdf}。在频域建模方面，FITS、FreTS 和 FilterNet 等工作利用快速傅里叶变换或可学习频率滤波器提取周期性信息，在保持轻量化的同时取得较强的时序表示能力\cite{fits,frets,filternet}。此外，真实时间序列通常存在分布漂移和非平稳性，AdaRNN、Dish-TS、RevIN 和 Non-stationary Transformers 等方法分别通过域适应、分布调整、可逆实例归一化和去平稳化缓解训练与测试分布不一致的问题\cite{adarnn,dishts,revin,ns-trans}。

时间戳是多元时间序列中区别于普通序列数据的重要信息。部分方法直接忽略时间戳\cite{patchtst,fits}；Autoformer 和 TimesNet 等模型则将历史时间戳编码为位置嵌入或辅助特征\cite{autoformer,timesnet}；iTransformer 进一步把时间戳作为独立令牌与观测序列共同输入注意力层\cite{itransformer}。近年来，研究者开始更显式地利用时间戳。GLAFF 通过自注意力分别建模历史和未来时间戳，并生成全局—局部信息的自适应权重\cite{glaff}；TimeLinear 通过编码器—解码器结构提取时间戳嵌入，再与数据嵌入加权融合；TimeSter 则系统讨论了时间相关特征对时序预测的影响\cite{timester}。然而，这些方法仍然或者只关注历史时间戳，或者将历史和未来时间戳分别处理，较少显式刻画二者之间的双向依赖关系。未来时间戳所携带的工作日、节假日、季节等先验上下文，对于修正仅由历史数据得到的预测结果具有重要意义。

\subsection{交通流预测中的时空图建模方法}

交通流预测是时空序列预测中最具代表性的任务之一，其核心是根据历史交通观测预测未来若干时刻的交通速度、流量或占有率\cite{traffic-stf-survey,survey01}。由于交通传感器或路段之间具有天然的图结构，图神经网络成为该类任务的主流建模工具。STGCN 将图卷积与时间卷积结合，首次系统地将时空图建模引入交通预测\cite{stgcn}；DCRNN 利用扩散卷积和编码器—解码器结构模拟交通信息的扩散过程\cite{dcrnn}；Graph WaveNet 通过学习自适应邻接矩阵，在无须给定显式图结构的条件下捕捉节点间依赖\cite{gwnet}。此后，ASTGCN、MTGNN、AGCRN 和 DGCRN 等工作分别从注意力时空卷积、图结构学习、自适应图生成和动态图建模等角度提升交通预测能力\cite{astgcn,mtgnn,agcrn,dgcrn}。STID 和 STAEformer 则进一步说明，合理的时空身份嵌入和自适应嵌入机制能够使简单 Transformer 在交通预测上取得较强性能\cite{stid,staeformer}。

真实交通数据在不同空间尺度上往往呈现不同的分布和依赖模式。为刻画这种空间非平稳性，研究者提出了多种空间多尺度方法。HGCN 基于距离邻接矩阵进行谱聚类，构造层次化空间图，并在各尺度上分别应用时空图神经网络\cite{hgcn}；HSTGL 分别从交通连通性和地理相关性出发学习微观图与宏观图\cite{hstgl}；STPGNN 通过节点特征相似度识别关键节点并构造新的邻接矩阵，在原始图与关键节点图两个尺度上提取特征\cite{stpgnn}；PyGNN 和 MAGNN 则从图采样、图过滤和多尺度自适应图学习等角度增强空间多尺度表示能力\cite{pygnn,magnn}。这些方法表明，将空间图划分为不同尺度进行分层建模，有助于同时捕捉局部路网结构和区域级交通模式。然而，现有方法大多依赖距离、聚类或节点相似性等先验知识构造空间尺度，并且通常只考虑原图与一个额外子图两个尺度，难以根据当前交通状态自适应确定尺度数量并进行灵活的跨尺度融合。

另一方面，交通数据中的静态路网结构和动态时空依赖通常同时存在。静态结构反映由道路连接决定的空间拓扑关系，动态依赖则反映随时间变化的交通传播和远程交互。D2STGNN 通过解耦静态图卷积与动态图卷积分别建模两类关系\cite{d2stgnn}；TESTAM 采用全激活的混合专家结构，分别用三个专家提取时间特征、静态图特征和动态图特征\cite{testam}；ST-MoE 则将时空混合专家用于交通预测中的去偏建模\cite{st-moe}。虽然混合专家机制为交通预测中的异质性建模提供了思路，但现有方法仍存在两个主要问题：一是部分方法在全量专家完成计算后才进行路由，未能实现真正的稀疏激活，限制了效率与可扩展性；二是硬专家选择策略难以在同一时间状态下联合建模静态空间结构和动态时间演化。

\subsection{多尺度建模与混合专家系统}

多尺度建模已在时序预测和交通预测中展现出明显优势。在时间维度上，不同粒度的时间序列分别携带趋势、季节和局部波动等信息，分层或并行处理多尺度表示有助于缓解单一尺度带来的信息损失\cite{pyraformer,scinet,pathformer,pdf}。在空间维度上，多尺度图表示能够刻画传感器网络从细粒度局部结构到粗粒度区域结构的非平稳模式\cite{hgcn,pygnn,magnn}。然而，简单地让所有尺度共享同一组参数，或仅依赖人工设定的尺度数量，难以适应输入分布的变化。尤其是多元时间序列中不同通道可能服从不同的时间分布，交通网络中不同区域也可能具有不同的空间模式，需要更加细粒度的参数分配和自适应融合策略。

混合专家系统为上述问题提供了一种自然的解决途径。经典稀疏门控 MoE 通过 Top-$k$ 路由在大量专家中只激活少量专家，从而在提高模型容量的同时控制计算开销\cite{moe,moe-survey}。SMoEs 将专家机制引入空间维度，以处理空间输入中的局部差异性\cite{smoes}；ST-MoE 和 TESTAM 则将 MoE 应用于交通预测，以应对时空分布的非平稳性\cite{st-moe,testam}。这些工作初步验证了专家路由在复杂时空建模中的有效性，但多尺度表示与混合专家机制之间仍缺乏系统性的结合。例如，如何依据时间分布为不同通道选择多尺度时序专家，如何依据空间结构为不同尺度图选择专门化的图专家，以及如何在稀疏激活条件下联合建模静态图与动态时空依赖，仍是尚待进一步研究的问题。本文正是围绕这些不足，分别从时间域和空间域构建多尺度建模与混合专家系统相结合的两类预测框架。

\section{本文的主要研究内容与贡献}

本文的核心研究问题可以概括为：在时空序列预测中，如何通过多尺度建模与混合专家系统，自适应地刻画数据在时间、变量和空间维度上的非平稳性与异质性，从而提升模型的预测精度和鲁棒性。围绕该问题，本文分别在一般多元时间序列预测和交通流预测两个代表性任务上开展研究，主要工作如下。

第一项工作针对一般多元时间序列预测任务，提出一种融合多尺度时序建模与混合专家系统的预测框架。现有方法往往忽视历史时间戳与未来时间戳之间的双向依赖关系，并且在多尺度时序建模时通常对所有通道共享参数，难以捕捉不同通道的数据分布差异。针对上述问题，该方法设计了时间戳双向交叉注意力模块（Timestamp Dual Cross-Attention Module，TDCAM），通过对称交叉注意力显式建模历史区间与未来区间之间的双向时间上下文，并将其作为约束融入预测过程；同时设计了多尺度特征提取混合专家模块（Multi-Scale Feature Extract Mixture-of-Experts，MSMoE），采用稀疏门控路由将不同通道的时序数据动态分配到若干专家中，每个专家对输入进行多尺度下采样、趋势—季节分解和跨尺度融合，从而实现对不同时间分布和多尺度模式的细粒度建模。此外，该方法还利用频域中的可学习马氏矩阵提取变量间相关关系，并引入辅助损失平衡专家负载。实验结果表明，该方法在 10 个真实世界数据集上取得了具有竞争力的预测性能，验证了双向时间戳建模和多尺度混合专家设计的有效性。

第二项工作针对时空交通预测任务，提出一种融合空间多尺度建模与时空静动态混合专家的预测框架。真实交通数据在不同空间尺度上呈现显著分布差异，同时存在由交通状态演化引起的时间非平稳性，且静态路网结构与动态时空依赖相互耦合。针对这些问题，该方法设计了空间多尺度混合专家模块（Spatial Multi-Scale Mixture-of-Experts，SMS-MoE），通过图池化以数据驱动方式构造多个空间尺度的子图，并利用尺度专家和自适应尺度融合机制提取尺度特异性特征；同时设计了静态—动态双混合专家模块（Static-Dynamic Dual Mixture-of-Experts，SD-MoE），其中静态图 MoE 使用轻量图卷积专家捕捉显式拓扑结构，动态图 MoE 使用节点级全注意力专家挖掘隐式时空依赖，并通过稀疏激活实现高效联合建模。最后，模型通过门控融合层将空间多尺度特征与静动态时空特征进行自适应融合。在 6 个公开交通数据集上的实验结果表明，该方法能够一致地取得先进预测性能，验证了自适应空间多尺度建模和静动态双专家学习的有效性。

从整体研究脉络来看，第一项工作侧重从时间、时间戳和通道分布角度解决一般多元时间序列中的非平稳与多尺度问题；第二项工作则将多尺度思想扩展到空间图结构，并进一步处理静态图与动态依赖共存的交通预测问题。二者共同构成了一套较为完整的“多尺度建模 + 混合专家”时空序列预测研究框架，分别从时间域和空间域验证了所提建模思想的通用性与有效性。

\section{本文章节组织}

本文围绕多尺度建模与混合专家系统在时空序列预测中的应用展开研究，全文共分为五章，各章内容安排如下。

第1章为绪论。本章介绍了时空序列预测的研究背景与意义，梳理了多元时间序列预测、交通流预测以及多尺度建模与混合专家系统等方面的国内外研究现状，分析了现有方法存在的不足，并概述了本文的主要研究内容与贡献。

第2章为相关理论基础。本章介绍时空序列预测问题定义、基于深度学习的多元时间序列预测方法、交通流预测中的图神经网络方法以及混合专家系统，为第3章和第4章的研究工作提供理论基础。

第3章为多尺度时序建模与混合专家系统的时间序列预测。本章面向一般多元时间序列预测任务，依次阐述问题定义、多尺度时序混合专家与双向时间戳依赖建模框架、时间戳双向交叉注意力模块、多尺度特征提取混合专家模块、频域变量相关建模方法以及整体损失函数，并给出实验设置、与主流方法的对比结果、消融实验和参数分析。

第4章为多尺度空间建模与混合专家系统的时空交通预测。本章面向时空交通预测任务，阐述数据驱动的空间多尺度图构造、空间多尺度混合专家模块、静态—动态双混合专家模块、特征融合与预测输出，并在多个交通预测基准数据集上验证模型性能及各模块的有效性。

第5章为总结与展望。本章对全文工作进行总结，归纳两项工作的研究结论和主要贡献，并讨论当前方法存在的不足以及未来可能的研究方向。
